\documentclass[letterpaper]{article} %Comienza el preámbulo
\usepackage{graphics}
\usepackage{graphicx}
\usepackage{amsmath}
\usepackage{amsfonts}
\usepackage[spanish]{babel}
\decimalpoint
\linespread{1.}
\usepackage[cm]{fullpage}
\usepackage{multicol}
\usepackage[utf8]{inputenc}
\usepackage{array}
\usepackage{enumerate,mdwlist}
\usepackage{framed}
\newcolumntype{L}[1]{>{\raggedright\let\newline\\\arraybackslash\hspace{0pt}}m{#1}}

\title{Universidad  \\Tercer Examen Parcial, Cálculo diferencial.  2016 - S1}


\newcommand{\grafica}[3]{\resizebox{#1}{#2}{\includegraphics{#3}}\\}
\newcommand{\figura}[2]{\resizebox{#1}{!}{\includegraphics{#2}}\\}
\newcommand{\beq}{\begin{equation}}
\newcommand{\eeq}{\end{equation}}
\newcommand{\bea}{\begin{eqnarray*}}
\newcommand{\eea}{\end{eqnarray*}}
\newcommand\printvalue{1}
\newcommand\printlabel{1}
\newcommand{\reviso}[1]{\ifnum\printlabel>0\textbf{Revisó:} #1 \\ \fi}
\newcommand{\anotaciones}[1]{\ifnum\printlabel>0\textbf{Anotaciones:} #1 \\ \fi}
\newcommand{\etiqueta}[1]{\ifnum\printlabel>0\textbf{#1} \\\fi}
\newcommand{\autor}[1]{\ifnum\printlabel>0 \textbf{Autor:} #1 \\ \fi }
\newcommand{\respuesta}[1]{\ifnum\printlabel>0 \textbf{Respuesta:}#1 \fi}
\newcommand{\temas}[1]{\ifnum\printlabel>0 \textbf{Temas:} #1 \\\fi}
\newcommand{\valor}[1]{\ifnum\printvalue>0\textit{(Calificación: +#1)}\fi}


\begin{document}
\date{}
\maketitle

Nombre: \verb+________________________________+ Grupo \verb+___________________+ Fecha: \verb+_________________+ \\

\begin{multicols}{2}
  \textbf{Compromiso ético:} Al entregar éste parcial, usted garantiza
  que ha intentado resolverlo con su mejor esfuerzo; y que las
  respuestas son el resultado únicamente de su trabajo y conocimiento
  individuales.

  \textbf{Instrucciones:} Guarde sus objetos personales. Tenga en su
  escritorio los útiles necesarios para resolver el parcial. No ponga
  maletas, bolsos y chaquetas en la mesa de trabajo. No use
  audífonos. Guarde su teléfono celular, sin sonido.

  Las preguntas $I$ y $II$ son de selección múltiple con
  única respuesta.
La calificación de cada una es $+0.7$.
Marque una
  única respuesta; en caso de que haya confusión en cuanto a cuál es
  su respuesta, se le calificará el punto como errado. No se califica
  procedimiento.

Las preguntas $III$, $IV$ y $V$ son de respuesta abierta. Sí se califica procedimiento. Tienen una calificación de $+1.2$ cada una.

% A B C D
% E F G H
% I J K L
% M N O P

%Tiene un peso de $0.5$:\\

\hspace{1cm}
    \begin{enumerate}[I]
\item Primera Pregunta (incluir el texto aquí)
\item Segunda Pregunta
\item ...
\end{enumerate}
\end{multicols}

\end{document}
